% GENERATED FILE. Edit canonical lessons, then run npm run generate:books.
% canonical-source-hash: fnv1a64:461eb99889d86bea
% canonical-lessons: RU-C502-glina, RU-C502-plastilin, RU-C502-gonchar, RU-C502-masterit, RU-C502-podelka

\chapter{Clay and Crafts}
\label{ch:ru-clay-and-crafts}
\hlchaptermodality{502}

\begin{chapteropening}
\textbf{By the end of this chapter:} \emph{I can talk about pottery, modelling and handmade things.}

Complete the last lesson of chapter 502: I can talk about pottery, modelling and handmade things.
\end{chapteropening}

\section[\texorpdfstring{\ru{глина} (glína)}{глина (glína)}]{\ru{глина} (glína) — clay}

\label{lesson:RU-C502-glina}



\begin{quote}
Before the new one: say the Russian for a sewing pattern, then the Russian for to hem.
\end{quote}

\subsection*{You'll want to know: \ru{глина}}
\textbf{\ru{глина}} (\emph{glína}) — \textquotedblleft{}clay\textquotedblright{}.

\begin{quote}
\ru{Из} \ru{глины} \ru{делают} \ru{кувшины} \ru{и} \ru{чашки}.
\end{quote}

\begin{quote}
\emph{Iz glíny délayut kuvshíny i cháshki.}
\end{quote}

\begin{quote}
Jugs and cups are made from clay.
\end{quote}

Feminine.

\begin{cousinweb}
\emph{\ru{глина}} is inherited from Proto-Slavic \emph{glina}, which EtymDB traces to PIE \emph{gleh\textsubscript{1}i-n-}.
\end{cousinweb}

\begin{grammarlens}[title={What you've built}]
One more word for telling, explaining and describing at B1.
\end{grammarlens}

\subsection*{Your turn}
\begin{itemize}
\raggedright
  \item \emph{Say it:} \emph{glína}
  \item \emph{Say it:} \emph{glína}, once more
  \item \emph{Say it:} a sentence of your own with \emph{glína}
  \item \emph{Recall:} say the Russian for a sewing pattern, then the Russian for to hem, then say \emph{glína} again
\end{itemize}

\begin{tcolorbox}[breakable,colback=teal!4,colframe=teal!35!black,title={Before you move on}]
What does \ru{глина} mean? (\textquotedblleft{}Clay\textquotedblright{}.) Say it once more.
\end{tcolorbox}
\section[\texorpdfstring{\ru{пластилин} (plastilín)}{пластилин (plastilín)}]{\ru{пластилин} (plastilín) — modelling clay}

\label{lesson:RU-C502-plastilin}



\begin{quote}
Before the new one: say the Russian for to hem, then the Russian for clay.
\end{quote}

\subsection*{You'll want to know: \ru{пластилин}}
\textbf{\ru{пластилин}} (\emph{plastilín}) — \textquotedblleft{}modelling clay\textquotedblright{}.

\begin{quote}
\ru{Дети} \ru{лепят} \ru{из} \ru{пластилина} \ru{зверей}.
\end{quote}

\begin{quote}
\emph{Déti lépyat iz plastilína zveréy.}
\end{quote}

\begin{quote}
The children are making animals out of modelling clay.
\end{quote}

Masculine.

\begin{cousinweb}
\emph{\ru{пластилин}} is usually traced to Italian \emph{plastilina}, built on Greek \emph{plastós} 'moulded'; EtymDB has no entry.
\end{cousinweb}

\begin{grammarlens}[title={What you've built}]
One more word for telling, explaining and describing at B1.
\end{grammarlens}

\subsection*{Your turn}
\begin{itemize}
\raggedright
  \item \emph{Say it:} \emph{plastilín}
  \item \emph{Say it:} \emph{plastilín}, once more
  \item \emph{Say it:} a sentence of your own with \emph{plastilín}
  \item \emph{Recall:} say the Russian for to hem, then the Russian for clay, then say \emph{plastilín} again
\end{itemize}

\begin{tcolorbox}[breakable,colback=teal!4,colframe=teal!35!black,title={Before you move on}]
What does \ru{пластилин} mean? (\textquotedblleft{}Modelling clay\textquotedblright{}.) Say it once more.
\end{tcolorbox}
\section[\texorpdfstring{\ru{гончар} (gonchár)}{гончар (gonchár)}]{\ru{гончар} (gonchár) — a potter}

\label{lesson:RU-C502-gonchar}



\begin{quote}
Before the new one: say the Russian for clay, then the Russian for modelling clay.
\end{quote}

\subsection*{You'll want to know: \ru{гончар}}
\textbf{\ru{гончар}} (\emph{gonchár}) — \textquotedblleft{}a potter\textquotedblright{}.

\begin{quote}
\ru{Гончар} \ru{делает} \ru{красивые} \ru{тарелки}.
\end{quote}

\begin{quote}
\emph{Gonchár délayet krasívyye tarélki.}
\end{quote}

\begin{quote}
The potter makes beautiful plates.
\end{quote}

Masculine; genitive \ru{гончара}.

\begin{cousinweb}
\emph{\ru{гончар}} is inherited from Old East Slavic \emph{gorn\ru{ь}čar\ru{ь}}, derived from \emph{g\ru{ъ}rn\ru{ь}c\ru{ь}} 'pot' (EtymDB).
\end{cousinweb}

\begin{grammarlens}[title={What you've built}]
One more word for telling, explaining and describing at B1.
\end{grammarlens}

\subsection*{Your turn}
\begin{itemize}
\raggedright
  \item \emph{Say it:} \emph{gonchár}
  \item \emph{Say it:} \emph{gonchár}, once more
  \item \emph{Say it:} a sentence of your own with \emph{gonchár}
  \item \emph{Recall:} say the Russian for clay, then the Russian for modelling clay, then say \emph{gonchár} again
\end{itemize}

\begin{tcolorbox}[breakable,colback=teal!4,colframe=teal!35!black,title={Before you move on}]
What does \ru{гончар} mean? (\textquotedblleft{}A potter\textquotedblright{}.) Say it once more.
\end{tcolorbox}
\section[\texorpdfstring{\ru{мастерить} (masterít)}{мастерить (masterít)}]{\ru{мастерить} (masterít) — to make things by hand}

\label{lesson:RU-C502-masterit}



\begin{quote}
Before the new one: say the Russian for modelling clay, then the Russian for a potter.
\end{quote}

\subsection*{You'll want to know: \ru{мастерить}}
\textbf{\ru{мастерить}} (\emph{masterít}) — \textquotedblleft{}to make things by hand\textquotedblright{}.

\begin{quote}
\ru{Папа} \ru{любит} \ru{мастерить} \ru{в} \ru{гараже}.
\end{quote}

\begin{quote}
\emph{Pápa lyúbit masterít v garazhé.}
\end{quote}

\begin{quote}
Dad likes making things in the garage.
\end{quote}

Imperfective; perfective \ru{смастерить}.

\begin{cousinweb}
\emph{\ru{мастерить}} is built on \emph{\ru{мастер}} 'master, craftsman', which EtymDB traces through Old French \emph{maistre} to Latin \emph{magister} 'master'; EtymDB lists the verb without a source.
\end{cousinweb}

\begin{grammarlens}[title={What you've built}]
One more word for telling, explaining and describing at B1.
\end{grammarlens}

\subsection*{Your turn}
\begin{itemize}
\raggedright
  \item \emph{Say it:} \emph{masterít}
  \item \emph{Say it:} \emph{masterít}, once more
  \item \emph{Say it:} a sentence of your own with \emph{masterít}
  \item \emph{Recall:} say the Russian for modelling clay, then the Russian for a potter, then say \emph{masterít} again
\end{itemize}

\begin{tcolorbox}[breakable,colback=teal!4,colframe=teal!35!black,title={Before you move on}]
What does \ru{мастерить} mean? (\textquotedblleft{}To make things by hand\textquotedblright{}.) Say it once more.
\end{tcolorbox}
\section[\texorpdfstring{\ru{поделка} (podélka)}{поделка (podélka)}]{\ru{поделка} (podélka) — a handmade craft item}

\label{lesson:RU-C502-podelka}



\begin{quote}
Before the new one: say the Russian for a potter, then the Russian for to make things by hand.
\end{quote}

\subsection*{You'll want to know: \ru{поделка}}
\textbf{\ru{поделка}} (\emph{podélka}) — \textquotedblleft{}a handmade craft item\textquotedblright{}.

\begin{quote}
\ru{Дети} \ru{показали} \ru{нам} \ru{свои} \ru{поделки} \ru{из} \ru{бумаги}.
\end{quote}

\begin{quote}
\emph{Déti pokazáli nam svoí podélki iz bumági.}
\end{quote}

\begin{quote}
The children showed us the things they had made out of paper.
\end{quote}

Feminine.

\begin{cousinweb}
\emph{\ru{поделка}} is built on \emph{\ru{поделать}}, from \emph{\ru{делать}} 'to do, to make'; EtymDB has no entry.
\end{cousinweb}

\begin{grammarlens}[title={What you've built}]
One more word for telling, explaining and describing at B1.
\end{grammarlens}

\subsection*{Your turn}
\begin{itemize}
\raggedright
  \item \emph{Say it:} \emph{podélka}
  \item \emph{Say it:} \emph{podélka}, once more
  \item \emph{Say it:} a sentence of your own with \emph{podélka}
  \item \emph{Recall:} say the Russian for a potter, then the Russian for to make things by hand, then say \emph{podélka} again
\end{itemize}

\begin{tcolorbox}[breakable,colback=teal!4,colframe=teal!35!black,title={Before you move on}]
What does \ru{поделка} mean? (\textquotedblleft{}A handmade craft item\textquotedblright{}.) Say it once more.
\end{tcolorbox}
